% ============================================================
%  BAB 3 – METODOLOGI PENELITIAN
%  Compile standalone:  pdflatex bab-3.tex
%  Compile as part of thesis: pdflatex main.tex
% ============================================================
\documentclass[../../thesis]{subfiles}
\usepackage{hyphenat}
\usepackage{iftex}

% ── Encoding, Language & Fonts ───────────────────────────────
\ifXeTeX
    \usepackage{polyglossia}
    \setmainlanguage{bahasai}
    \usepackage{fontspec}
    \setmainfont{TeX Gyre Termes}
    \setsansfont{TeX Gyre Heros}
    \setmonofont[Scale=0.88]{TeX Gyre Cursor}
\else
    \usepackage[utf8]{inputenc}
    \usepackage[T1]{fontenc}
    \usepackage[indonesian]{babel}
    \usepackage{mathptmx}
    \usepackage{helvet}
    \usepackage{courier}
\fi

% ── Page geometry (A4, UI thesis margins) ────────────────────
\usepackage[
  a4paper,
  top    = 2.54cm,
  bottom = 2.54cm,
  left   = 3.17cm,
  right  = 2.54cm
]{geometry}

% ── Line spacing & paragraph ─────────────────────────────────
\usepackage{setspace}
\onehalfspacing
\setlength{\parindent}{1.27cm}
\setlength{\parskip}{0pt}

% ── Microtype & justification ────────────────────────────────
\usepackage{microtype}
\sloppy

% ── Headings ─────────────────────────────────────────────────
\usepackage{titlesec}

\titleformat{\chapter}[block]
  {\centering\rmfamily\bfseries\large}{}
  {0pt}{\MakeUppercase}
\titlespacing*{\chapter}{0pt}{0pt}{1.5em}

\titleformat{\section}
  {\rmfamily\bfseries\normalsize}{\thesection}{1em}{}
\titlespacing*{\section}{0pt}{1.2em}{0.6em}

\titleformat{\subsection}
  {\rmfamily\bfseries\normalsize}{\thesubsection}{1em}{}
\titlespacing*{\subsection}{0pt}{1em}{0.5em}

\titleformat{\subsubsection}
  {\rmfamily\bfseries\normalsize}{\thesubsubsection}{1em}{}
\titlespacing*{\subsubsection}{0pt}{0.8em}{0.4em}

% ── Tables ───────────────────────────────────────────────────
\usepackage{longtable, booktabs, array, multirow}
\usepackage{calc}

% ── Math ─────────────────────────────────────────────────────
\usepackage{amsmath, amssymb}

% ── Graphics ─────────────────────────────────────────────────
\usepackage{graphicx}
\makeatletter
\def\maxwidth{\ifdim\Gin@nat@width>\linewidth\linewidth\else\Gin@nat@width\fi}
\makeatother
\setkeys{Gin}{width=\maxwidth,keepaspectratio}

% ── TikZ (ontology diagram) ──────────────────────────────────
\usepackage{tikz}
\usetikzlibrary{shapes.geometric, arrows.meta, positioning, calc, fit, backgrounds}

% ── Extended tables ──────────────────────────────────────────
\usepackage{tabularx}

% ── Lists ────────────────────────────────────────────────────
\usepackage{enumitem}
\setlist[enumerate,1]{
  label      = \arabic*.,
  leftmargin = 1.27cm,
  itemsep    = 4pt,
  parsep     = 0pt,
  topsep     = 4pt
}
\setlist[enumerate,2]{
  label      = \alph*.,
  leftmargin = 1.27cm,
  itemsep    = 3pt,
  parsep     = 0pt,
  topsep     = 2pt
}

% ── Hyperlinks ───────────────────────────────────────────────
\usepackage[hidelinks]{hyperref}

% Bibliography setup is inherited from main.tex (biblatex + references.bib)
% when this file is compiled standalone via the subfiles class.

% ============================================================
\begin{document}

% When standalone: set counter so \chapter{} renders as "BAB 3".
% When in main.tex: main.tex sets the counter before \subfile{bab-3}.
\ifSubfilesClassLoaded{}{\setcounter{chapter}{2}}

\chapter{Metodologi Penelitian}

Tinjauan pada Bab~2 menyimpulkan bahwa \textit{Knowledge Graph Completion}
(KGC) berbantuan \textit{Large Language Model} (LLM) merupakan pendekatan yang
menjanjikan, tetapi belum teruji untuk konteks kurikulum sains lintas-disiplin
berbahasa Indonesia. Menguji pendekatan tersebut secara sistematis menuntut
rancangan yang menjembatani gagasan konseptual dengan artefak yang dapat
dibangun dan dinilai. Bab ini menetapkan rancangan itu dalam paradigma
\textit{Design Science Research}, menguraikan bagaimana data kurikulum diperoleh
dan diproses, bagaimana skema ontologi serta \textit{pipeline} empat-tahap
dirancang, dan bagaimana kualitas graf yang dihasilkan diukur, baik secara
struktural maupun melalui validasi pakar, agar klaim efektivitas pendekatan
dapat dipertanggungjawabkan.

% ============================================================
\section{Desain Penelitian}
% ============================================================

Penelitian ini menggunakan pendekatan \textit{Design Science Research}
(DSR) yang difokuskan pada pengembangan purwarupa sistem KGC berbasis
LLM. DSR dipilih karena penelitian ini secara spesifik bertujuan
menghasilkan artefak inovatif, representasi pengetahuan kurikulum yang
bersifat lintas-disiplin, sebagai respons terhadap permasalahan
fragmentasi pembelajaran sains pada jenjang SMA.

Menurut \textcite{hevner_design_2004}, DSR adalah paradigma pemecahan masalah
yang berupaya meningkatkan kapabilitas manusia dan organisasi dengan
menciptakan artefak baru, direpresentasikan dalam bentuk konstruk,
model, metode, dan instansiasi. Artefak yang dikembangkan dalam
penelitian ini meliputi:

\begin{enumerate}
  \item \textbf{Model (Skema Ontologi):} Skema ontologi kurikulum sains yang
        dinyatakan dalam \textit{Web Ontology Language} (OWL), mencakup definisi
        kelas, relasi objek (\textit{object property}), dan properti data
        (\textit{datatype property}). Skema ini menjadi kerangka semantik yang
        menstrukturkan KG dan dirancang pada Subbab~\ref{subsec:skema-ontologi}.
  \item \textbf{Metode (Algoritma dan \textit{Pipeline}):} Algoritma
        ekstraksi pengetahuan terstruktur menggunakan LLM dengan
        \textit{output} berbasis ontologi, serta \textit{pipeline} KG
        \textit{Construction} dan \textit{Completion} yang diuraikan pada
        Subbab~\ref{subsec:kerangka-konsep}.
  \item \textbf{Instansiasi (Purwarupa Sistem):} \textit{Knowledge Graph}
        kurikulum sains yang terinstansiasi pada basis data Neo4j beserta
        visualisasinya, sebagai purwarupa struktur data yang dapat ditelusuri
        dan diakses secara komputasional.
\end{enumerate}

Secara metodologis, penelitian ini dikategorikan sebagai penelitian
terapan (\textit{applied research}) sebagaimana dirangkum pada
Tabel~\ref{tab:karakteristik}.

\begin{table}[htbp]
\centering
\caption{Karakteristik Metodologi Penelitian}
\label{tab:karakteristik}
\begin{tabular}{@{}ll@{}}
  \toprule
  \textbf{Karakteristik} & \textbf{Deskripsi} \\
  \midrule
  Jenis Penelitian & Terapan (\textit{Applied Research}) \\
  Pendekatan       & \textit{Design Science Research} (DSR) \\
  Metode           & Eksperimental dengan validasi kualitatif pakar \\
  Output           & Artefak sistem (Skema Ontologi, Metode KGC, dan Purwarupa Visualisasi) \\
  Paradigma        & Pragmatisme (berfokus pada utilitas dan pemecahan masalah) \\
\end{tabular}
\end{table}

Adopsi DSR dalam penelitian ini mengikuti kerangka kerja
\textcite{peffers_design_2007}, yang meliputi enam tahapan berikut.

\begin{enumerate}
  \item \textbf{Identifikasi Masalah:} Fragmentasi kurikulum sains
        lintas-disiplin (Fisika, Kimia, Biologi) menyebabkan siswa
        kesulitan melihat keterhubungan konsep.
  \item \textbf{Definisi Tujuan:} Mengembangkan sistem yang dapat
        menemukan relasi implisit antarkonsep lintas mata pelajaran.
  \item \textbf{Desain dan Pengembangan:} Membangun \textit{pipeline}
        ekstraksi--konstruksi--penyempurnaan \textit{knowledge graph}.
  \item \textbf{Demonstrasi:} Menerapkan artefak pada buku teks resmi
        Kemendikbudristek Fisika, Kimia, dan Biologi Kelas~XII.
  \item \textbf{Evaluasi:} Pengukuran metrik struktural (AD,
        \textit{Modularity}), metrik ekstraksi, dan validasi pakar.
  \item \textbf{Komunikasi:} Dokumentasi dalam bentuk skripsi dan
        publikasi ilmiah.
\end{enumerate}

% ============================================================
\section{Lokasi dan Waktu Penelitian}
% ============================================================

\subsection{Lokasi Penelitian}

Penelitian ini dilaksanakan secara daring tanpa melibatkan studi
lapangan. Seluruh proses ekstraksi, konstruksi graf, dan pengumpulan
data validasi dilakukan melalui platform digital.

\subsection{Waktu Penelitian}

Penelitian berlangsung selama enam bulan, dari Januari hingga Juni 2026.
Rincian tahapan disajikan pada Tabel~\ref{tab:waktu}.

\begin{table}[htbp]
\centering
\caption{Rincian Waktu Penelitian}
\label{tab:waktu}
\begin{tabular}{@{}lll@{}}
  \toprule
  \textbf{Fase} & \textbf{Kegiatan} & \textbf{Periode} \\
  \midrule
  1 & Studi literatur dan desain sistem
    & Januari -- Februari 2026 \\
  2 & Implementasi \textit{pipeline} ekstraksi
    & Februari -- Maret 2026 \\
  3 & Konstruksi \textit{Knowledge Graph}
    & \multirow{2}{*}{Maret -- April 2026} \\
  4 & Evaluasi dan validasi pakar \#1 & \\
  5 & \textit{Knowledge Graph Completion}
    & April 2026 \\
  6 & Evaluasi dan validasi pakar \#2
    & Mei 2026 \\
  7 & Analisis hasil dan penulisan
    & Mei -- Juni 2026 \\
\end{tabular}
\end{table}

% ============================================================
\section{Sumber Data dan Korpus Penelitian}
% ============================================================

Penelitian ini tidak menggunakan populasi manusia sebagai objek
eksperimen utama, melainkan dua jenis data yang saling melengkapi: data
tekstual sebagai korpus ekstraksi dan data penilaian pakar sebagai dasar
evaluasi model.

\subsection{Data Primer (Korpus Ekstraksi)}\label{subsec:data-primer}

Data primer yang digunakan dibatasi pada materi sains tingkat SMA
Fase~F (Kelas~XII). Pemilihan dokumen dilakukan secara
\textit{purposive sampling} berdasarkan relevansinya dengan Kurikulum
Merdeka. Korpus penelitian terdiri dari tiga buku teks siswa
(\textit{e-book}) terbitan Kemendikbudristek untuk mata pelajaran Fisika,
Kimia, dan Biologi Fase~F, yang diperoleh dari repositori resmi
\url{https://buku.kemdikbud.go.id}. Pembatasan pada buku teks resmi
memastikan bahwa entitas materi dan relasi dasar yang diekstrak sistem
benar-benar mencerminkan standar materi yang diajarkan secara nasional.
Rincian bibliografis ketiga buku siswa beserta Buku Panduan Guru
pasangannya (judul, penulis, dan ISBN) disajikan pada \ref{lamp:korpus}.

Secara kuantitatif, korpus mencakup 752 halaman Buku Siswa dan 736 halaman Buku
Panduan Guru pasangannya, dengan struktur kurikulum acuan sebanyak 17 bab dan 77
subbab yang dirujuk dari daftar isi Buku Siswa. Untuk diproses oleh
\textit{pipeline}, teks dipecah menjadi \textit{chunk} berukuran 800 token
dengan tumpang-tindih (\textit{overlap}) 200 token: 782 \textit{chunk} dari Buku
Siswa ($\approx$139.200 kata) dan 747 \textit{chunk} dari Buku Guru, sehingga
total korpus berjumlah 1.529 \textit{chunk}. Jumlah kata Buku Guru tidak
terhitung penuh karena pada Buku Guru Biologi setiap halaman tersimpan sebagai
gambar dan bukan teks, sehingga jumlah katanya tidak bisa dihitung otomatis
(lihat catatan Tabel~\ref{tab:ukuran-korpus}). Estimasi dari dua buku guru
lainnya sekitar 79.300 kata.
Rincian ukuran korpus per mata pelajaran disajikan pada
Tabel~\ref{tab:ukuran-korpus}.

\begin{table}[htbp]
\centering
\caption{Ukuran Korpus Ekstraksi per Mata Pelajaran}
\label{tab:ukuran-korpus}
\small
\setlength{\tabcolsep}{5pt}
\begin{tabular}{@{}lrrrrrrrr@{}}
  \toprule
  \multirow{2}{*}{\textbf{Mata Pelajaran}}
    & \multicolumn{2}{c}{\textbf{Halaman}}
    & \multirow{2}{*}{\textbf{Bab}} & \multirow{2}{*}{\textbf{Subbab}}
    & \multicolumn{2}{c}{\textbf{\textit{Chunk}}}
    & \multicolumn{2}{c}{\textbf{Kata ($\approx$)}} \\
  \cmidrule(lr){2-3} \cmidrule(lr){6-7} \cmidrule(lr){8-9}
    & \textbf{Siswa} & \textbf{Guru} & & & \textbf{Siswa} & \textbf{Guru}
    & \textbf{Siswa} & \textbf{Guru} \\
  \midrule
  Fisika  & 248 & 256 & 9 & 39 & 249 & 264 & 45.300 & 48.700 \\
  Kimia   & 224 & 216 & 4 & 22 & 234 & 219 & 36.900 & 30.600 \\
  Biologi & 280 & 264 & 4 & 16 & 299 & 264 & 57.000 & --\textsuperscript{*} \\
  \midrule
  \textbf{Total} & \textbf{752} & \textbf{736} & \textbf{17} & \textbf{77}
    & \textbf{782} & \textbf{747} & \textbf{139.200} & \textbf{79.300\textsuperscript{+}} \\
  \bottomrule
\end{tabular}

\vspace{4pt}
\begin{minipage}{\textwidth}
{\footnotesize\setstretch{1.1}\noindent
\textsuperscript{*}~Pada berkas Buku Guru Biologi
(\texttt{Biologi\_BG\_KLS\_XII.pdf}), setiap halaman tersimpan dalam bentuk
gambar dan bukan teks, sehingga jumlah katanya tidak bisa dihitung secara
otomatis. Karena itu, jumlah \textit{chunk}-nya (264) sama persis dengan jumlah
halaman dan tidak berisi teks.
\textsuperscript{+}~Total kata Buku Guru hanya dihitung dari Fisika dan Kimia.
Kolom Bab dan Subbab tidak dipisah Siswa/Guru karena mengikuti daftar isi Buku
Siswa.\par}
\end{minipage}
\end{table}

\subsection{Data Validasi Pakar}\label{subsec:data-validasi}

Data validasi diperoleh melalui penilaian kualitatif oleh panel pakar
terhadap \textit{triple} relasi (khususnya interkoneksi
lintas-disiplin) yang dihasilkan sistem KGC. Setiap \textit{triple}
dinilai menggunakan empat kategori penilaian yang telah didefinisikan pada
Subbab~\ref{subsec:validasi-manusia}, yaitu Benar (1,0), Sebagian Benar (0,5), Kurang Konteks
(0,25), dan Salah (0,0). Selain memberikan penilaian, pakar dapat
menambahkan komentar dan \textit{missing triples} yang dianggap esensial
namun tidak ditemukan sistem.

Pengumpulan data validasi dilakukan melalui aplikasi web KG Review App
yang dikembangkan khusus untuk penelitian ini (arsitektur teknis dibahas
pada Subbab~\ref{sec:instrumen-evaluasi}). Pada Fase~1, penilaian awal dari
dua pakar di KG Review App dan label konsensus akhir digunakan bersama:
label konsensus menjadi \textit{ground truth} untuk menghitung
Fuzzy Precision, Fuzzy Recall, dan F1-Score, sedangkan dua penilaian asli tersebut
digunakan untuk menguji tingkat kesepakatan antarpakar menggunakan
Gwet's AC1.

% ============================================================
\section{Subjek Evaluasi dan Panel Pakar}\label{subsec:subjek-evaluasi}
% ============================================================

Evaluasi performa sistem menggunakan pendekatan \textit{human
evaluation}. Subjek penelitian adalah panel pakar domain yang bertugas
sebagai pakar untuk memvalidasi akurasi ontologis dan
relevansi pedagogis relasi lintas-disiplin yang dihasilkan sistem.

Penelitian ini melibatkan enam pakar yang mewakili mata
pelajaran Fisika, Kimia, dan Biologi (dua pakar per mata
pelajaran); profil keilmuan masing-masing pakar disajikan pada
\ref{lamp:profil-pakar}. Jumlah dan kriteria pakar ditetapkan secara
\textit{purposive sampling} berdasarkan tiga pertimbangan metodologis
berikut.

\begin{enumerate}
  \item \textbf{Pengukuran Reliabilitas Multi-Rater:} Pada Fase~1, skema
        evaluasi mengharuskan satu \textit{triple} dalam-buku divalidasi
        oleh dua pakar independen, yang memungkinkan perhitungan
        \textit{inter-rater agreement} menggunakan Gwet's AC1. Adapun
        \textit{edge} lintas-mapel pada Fase~2 dinilai bersama melalui
        diskusi panel.

  \item \textbf{Beban Kognitif Penilaian:} Jumlah \textit{triple}
        relasi per mata pelajaran yang dihasilkan sistem melebihi 100
        item. Jumlah pakar dibatasi agar tidak menurunkan ketelitian
        akibat beban kognitif yang berlebihan selama proses anotasi.

  \item \textbf{Keketatan Kualifikasi:} Pakar harus
        merupakan mahasiswa aktif atau lulusan dengan latar belakang
        jurusan ilmu murni atau pendidikan yang linier (Biologi, Kimia,
        atau Fisika). Penggunaan mahasiswa dengan keahlian domain yang
        linier mengikuti preseden studi anotasi KG pendidikan yang
        mengutamakan kedalaman keilmuan spesifik dibanding pengalaman
        praktis umum \parencite{wei_kicgpt_2023}. Ketentuan ini secara alami
        membatasi ketersediaan populasi yang memenuhi kualifikasi.
\end{enumerate}

% ============================================================
\section{Kerangka Konsep}\label{subsec:kerangka-konsep}
% ============================================================

Kerangka konsep penelitian ini menggambarkan \textit{pipeline}
pemrosesan data dari dokumen kurikulum hingga \textit{Knowledge Graph
Completion}. Alur penelitian tidak hanya terdiri atas ekstraksi dan
\textit{completion}, tetapi juga mencakup perantara artefak JSON,
penayangan graf pada aplikasi validasi, konsensus pakar, pemuatan ulang
graf konsensus ke Neo4j, serta pengukuran metrik sebelum dan sesudah
\textit{completion}.

\subsection{Diagram Alur Penelitian}

\textit{Pipeline} penelitian mengikuti alur artefak sebagaimana
diilustrasikan pada Gambar~\ref{fig:pipeline}. Setiap tahap memiliki
masukan dan keluaran eksplisit agar perbedaan antara graf hasil
ekstraksi, graf konsensus, dan graf pasca-\textit{completion} tidak
tercampur.

\subfile{pipeline-kgc}

\noindent\textbf{Masukan dan keluaran tiap tahap:}

\begin{enumerate}
  \item \textbf{KG Construction.} Masukan tahap ini adalah PDF Buku Siswa
        dan Buku Guru Kelas XII. Prosesnya mencakup ekstraksi struktur
        hierarki dari daftar isi, ekstraksi terminologi glosarium dan
        materi pokok, pemotongan teks menjadi \textit{chunk}, ekstraksi
        konsep dan relasi dalam-buku dengan LLM yang
        \textit{schema-constrained}, serta penyusunan struktur graf.
        Keluarannya adalah berkas KG JSON awal per mata pelajaran. Ingest
        ke Neo4j pada titik ini bersifat opsional dan digunakan untuk
        inspeksi atau penayangan graf, bukan sebagai sumber konsensus
        final.

  \item \textbf{KG Review App dan review pakar.} KG JSON awal
        ditampilkan dalam KG Review App agar pakar dapat menilai relasi
        dalam-buku, memberi label kebenaran, dan mengusulkan
        \textit{missing triple}. Keluarannya adalah berkas penilaian
        pakar per mata pelajaran.

  \item \textbf{Final consensus.} Berkas KG JSON awal, penilaian pakar,
        dan konteks buku teks diproses pada \textit{notebook} konsensus.
        Ketidaksepakatan antarpakar diajudikasi, yaitu ditengahi untuk
        menetapkan satu label final ketika dua pakar memberi penilaian yang
        berbeda atau ketika hanya satu pakar mengusulkan \textit{missing
        triple}. Proses ini menggunakan \textit{LLM-as-a-judge} berbasis
        konteks sumber; rincian alur ajudikasi ditunjukkan pada
        Gambar~\ref{fig:llm-judge}. Hasil
        akhirnya disimpan sebagai \texttt{\{mapel\}.consensus.json} dan
        \textit{gold standard}. Graf konsensus ini diingest kembali ke
        Neo4j dengan membentuk hierarki struktural
        \texttt{Grade--Chapter--Subtopic--Concept} dan relasi semantik
        dalam-buku untuk menghitung metrik struktural. Adapun berkas penilaian
        awal pakar dari KG Review App dan \textit{gold standard} digunakan
        bersama untuk menghitung metrik validasi Fase~1: Gwet's AC1 dihitung
        dari pasangan penilaian pakar, sedangkan Fuzzy Precision, Fuzzy Recall, dan F1-Score
        dihitung dengan membandingkan label awal pakar terhadap label final
        pada \textit{gold standard}/konsensus.

  \item \textbf{KG Completion.} Masukan tahap ini adalah graf konsensus
        yang telah diingest ke Neo4j dari
        \texttt{\{mapel\}.consensus.json}. Graf tersebut telah memuat relasi
        dalam-buku hasil ajudikasi. Proses \textit{completion} membangun
        \textit{embedding} konsep pada \textit{node} Neo4j, mencari kandidat
        lintas-buku dengan ANN pada Neo4j Vector Index, menyaring kandidat agar
        hanya mencakup pasangan lintas-mata-pelajaran yang belum memiliki
        relasi tertipe, lalu mengklasifikasikan relasi dengan LLM dan
        menuliskannya kembali ke graf. Keluarannya adalah daftar \textit{edge}
        \texttt{LINTAS\_BUKU\_*} baru beserta \textit{confidence score}
        dan deskripsi singkat.

  \item \textbf{Evaluasi pasca-\textit{completion}.} \textit{Edge}
        lintas-buku hasil \textit{completion} terlebih dahulu dituliskan
        ke graf Neo4j pasca-KGC. Graf ini kemudian digunakan untuk
        menghitung metrik struktural pasca-KGC dan membandingkan kondisi
        pra-KGC dan pasca-KGC. Relasi tersebut juga dibahas melalui panel
        pakar lintas-disiplin untuk menilai validitas, ketepatan tipe,
        dan ketepatan arah relasi.
\end{enumerate}

\subfile{llm-judge}

\subsection{Skema Ontologi Knowledge Graph}
\label{subsec:skema-ontologi}

Ontologi dalam penelitian ini diposisikan menggunakan dua tier
representasi yang dibedakan secara metodologis. \textbf{Tier
ekstraksi} menghasilkan model konseptual \textit{deskriptif} atas isi
buku teks Kurikulum Merdeka (konsep, sub-bab, bab, beserta
keterkaitan semantik antar-konsep dalam satu buku), sejalan dengan
tradisi \textit{domain module} dalam konstruksi otomatis \textit{knowledge
graph} dari buku teks elektronik. \textbf{Tier completion} menghasilkan kandidat
keterkaitan lintas-buku (\texttt{LINTAS\_BUKU\_*}) sebagai inferensi
pedagogis dangkal di atas Tier ekstraksi: sistem mengusulkan jembatan
antar-disiplin yang dapat dipertimbangkan, tetapi tidak menetapkannya
sebagai jalur belajar.

Predikat ``dangkal'' bersifat operasional: sistem hanya menginferensi
keterkaitan \textit{single-hop} antar dua konsep di buku berbeda dan
tidak menyusun jalur belajar multi-langkah lintas-disiplin (batasan
inferensi pedagogis dirinci pada Subbab~\ref{sec:ruang-lingkup}).
Dengan demikian, relasi berarah pada Tier ekstraksi seperti
\texttt{PRASYARAT} dan \texttt{MEMPERSIAPKAN} dimaknai sebagai
ketergantungan konseptual yang terdokumentasikan dalam materi sumber
(keputusan editorial penyusun buku), sedangkan relasi lintas-buku pada Tier
completion seperti \texttt{LINTAS\_BUKU\_PRASYARAT\_UNTUK} dimaknai sebagai
sugesti per-pasangan yang dapat dipertimbangkan secara lokal, bukan anjuran
kronologis bagi siswa.

Skema \textit{knowledge graph} diformalkan sebagai ontologi OWL untuk
memastikan konsistensi struktur dan vokabulari relasi antartahap
\textit{pipeline} (ekstraksi, konstruksi, dan \textit{completion}).
Topologi keseluruhan mencakup relasi struktural, relasi keterurutan
penyajian materi dalam buku (\textit{NEXT\_CHAPTER}), dan relasi semantik
pada Tier ekstraksi, serta inferensi lintas-buku pada Tier completion.
Keseluruhan topologi ini diilustrasikan pada Gambar~\ref{fig:ontologi}.

\subfile{onthology}

% ============================================================
\section{Parameter Sistem dan Definisi Operasional}
% ============================================================

\subsection{Parameter Sistem}
\label{subsec:parameter-sistem}

Parameter penelitian ini adalah variabel-variabel \textit{pipeline} yang
dapat dimodifikasi peneliti dan memengaruhi karakteristik \textit{knowledge
graph} yang dihasilkan. Variabel-variabel tersebut dirangkum pada
Tabel~\ref{tab:parameter}.

\begin{table}[htbp]
\centering
\caption{Parameter Sistem \textit{Pipeline}}
\label{tab:parameter}
\begin{tabular}{@{}p{0.30\linewidth}p{0.35\linewidth}p{0.27\linewidth}@{}}
  \toprule
  \textbf{Parameter} & \textbf{Deskripsi} & \textbf{Nilai/Level} \\
  \midrule
  Model LLM Ekstraksi
    & Model bahasa besar untuk ekstraksi konsep dan relasi dalam-buku
    & {\scriptsize \texttt{gemini-2.5-flash-lite} (akses Mei 2026)} \\
  Model LLM Completion
    & Model bahasa besar untuk klasifikasi relasi lintas-disiplin pada tahap KGC
    & {\scriptsize \texttt{gemini-2.5-flash} (akses Mei 2026)} \\
  Model LLM-as-a-judge
    & Model penengah ketidaksepakatan pakar pada \textit{pipeline} validasi
    & {\scriptsize \texttt{gemini-3.1-pro-preview} (akses Mei 2026)} \\
  Strategi \textit{Prompt Engineering}
    & Variasi instruksi sistem untuk ekstraksi
    & \textit{Schema-constrained few-shot} + \textit{context injection}
      + \textit{strict output constraints} \\
  Model \textit{Embedding}
    & Model representasi vektor untuk perhitungan similaritas
    & \texttt{gemini-embedding-001} (Gemini Embedding) \\
  Threshold Similaritas
    & Ambang batas \textit{cosine similarity} untuk KGC lintas-buku
    & 0,70 \\
  Parameter \textit{Chunking}
    & Pemotongan dokumen sebelum ekstraksi agar konteks stabil
    & \texttt{chunk\_size} = 800, \texttt{chunk\_overlap} = 200 \\
  Sumber Data
    & Domain materi yang diproses \textit{pipeline}
    & Biologi, Fisika, Kimia Kelas~XII \\
\end{tabular}
\end{table}

\subsection{Metrik Evaluasi}\label{subsec:metrik-evaluasi}

Metrik evaluasi dalam penelitian ini mencakup dimensi topologi graf dan
dimensi validasi pakar, sebagaimana dirangkum pada
Tabel~\ref{tab:metrik}. Definisi operasional lengkap setiap metrik
struktural telah dijabarkan pada Subbab~\ref{subsec:metrik-topologi}
dan~\ref{subsec:metrik-ekstraksi}; definisi
metrik validasi pakar pada Subbab~\ref{subsec:validasi-manusia}.

\begin{table}[htbp]
\centering
\caption{Metrik Evaluasi Penelitian}
\label{tab:metrik}
\begin{tabular}{@{}p{0.35\linewidth}p{0.37\linewidth}p{0.20\linewidth}@{}}
  \toprule
  \textbf{Variabel} & \textbf{Deskripsi} & \textbf{Satuan} \\
  \midrule
  Jumlah Konsep Terekstrak
    & Volume konsep hasil ekstraksi per dokumen
    & Jumlah \textit{node} \\
  Jumlah Relasi Terprediksi
    & Volume relasi lintas-buku hasil KGC
    & Jumlah \textit{edge} \\
  \textit{Average Degree} (AD)
    & Tingkat keterhubungan rata-rata; dibandingkan pre- dan post-KGC
    & Rata-rata derajat ($\geq 0$) \\
  \textit{Graph Density}
    & Kepadatan relasi graf; dibandingkan pre- dan post-KGC
    & Nilai 0--1 \\
  \textit{Modularity}
    & Keterpisahan komunitas antarmapel; dibandingkan pre- dan post-KGC
    & Nilai $-0{,}5$ hingga $1$ \\
  \textit{Description Completeness} (DC)
    & Kelengkapan deskripsi entitas
    & Nilai 0--1 \\
  \textit{Empty SubKonsep Rate} (ESR)
    & Proporsi entitas induk tanpa turunan
    & Nilai 0--1 \\
  \textit{Typed Relation Ratio} (TRR)
    & Proporsi relasi bertipe spesifik dari kosakata ontologi
    & Nilai 0--1 \\
  \textit{Fuzzy Precision}
    & Ketepatan relasi berbasis skor kebenaran parsial (validasi pakar)
    & Nilai 0--1 \\
  \textit{Fuzzy Recall}
    & Cakupan relasi esensial berbasis skor kebenaran parsial (validasi pakar)
    & Nilai 0--1 \\
  F1-Score
    & Rata-rata harmonik \textit{Fuzzy Precision} dan \textit{Fuzzy Recall}
    & Nilai 0--1 \\
  Gwet's AC1
    & Kesepakatan antarpakar yang dikoreksi terhadap kebetulan
    & Nilai $-1$ hingga $1$ \\
\end{tabular}
\end{table}

\subsection{Instrumen dan Metrik Evaluasi Fase 2}\label{subsec:metrik-fase2}

Metrik \textit{Fuzzy Precision}, \textit{Fuzzy Recall}, F1-Score, dan Gwet's AC1 pada
Tabel~\ref{tab:metrik} berlaku untuk evaluasi Fase~1 (intra-mapel), yang menilai
setiap \textit{triple} dalam-buku pada skala empat kategori
(Subbab~\ref{subsec:validasi-manusia}). Evaluasi Fase~2 (inter-mapel)
menggunakan instrumen dan perlakuan metrik yang berbeda, karena objek dan
prosedur penilaiannya tidak sama.

\textbf{Instrumen.} Pada Fase~2, setiap relasi lintas-mapel $A \to B$ dinilai
panel pakar pada tiga aspek terpisah: (1) \textit{validitas}, yaitu apakah konsep
$A$ dan $B$ benar-benar berkaitan (Ya/Tidak); (2) \textit{ketepatan tipe}, yaitu
apakah label relasi yang diusulkan tepat (Benar/Salah); dan (3)
\textit{ketepatan arah}, yaitu apakah arah $A \to B$ sudah benar
(Benar/Terbalik/NA untuk relasi simetris). Relasi disajikan secara acak dan
\textit{score-blind}. Berbeda dengan Fase~1 yang melibatkan dua pakar
independen per \textit{triple}, Fase~2 menghasilkan satu jawaban konsensus per
relasi melalui diskusi panel.

\textbf{Metrik.} Dari ketiga aspek tersebut diturunkan tiga proporsi:
\textit{validitas relasi terprediksi}, yaitu proporsi kandidat relasi
lintas-mapel hasil sistem yang dinilai valid oleh panel, serta \textit{ketepatan
tipe} dan \textit{ketepatan arah} sebagai proporsi relasi dengan tipe dan arah
yang benar. Nilai validitas ini bersifat mirip \textit{precision} karena
penyebutnya adalah seluruh relasi yang diprediksi sistem, tetapi tidak identik
dengan \textit{Fuzzy Precision} Fase~1 karena tidak memakai skala parsial dan
tidak dibandingkan dengan dua penilaian pakar independen. Tiga hal membedakan
perlakuan metrik Fase~2 dari Fase~1:

\begin{itemize}
  \item \textbf{Tidak ada skala parsial.} Penilaian Fase~2 bersifat
        kategorikal per aspek, sehingga metrik berbasis himpunan \textit{fuzzy}
        tidak digunakan.
  \item \textbf{\textit{Fuzzy Recall} dan F1 tidak dihitung.} Tidak tersedia himpunan
        rujukan menyeluruh atas seluruh relasi lintas-mapel yang
        \emph{seharusnya} ada (sebagai penyebut \textit{Fuzzy Recall}); himpunan yang
        dievaluasi adalah seluruh kandidat yang dihasilkan sistem, bukan sampel
        dari populasi relasi benar yang telah diketahui.
  \item \textbf{Gwet's AC1 tidak berlaku.} Karena setiap relasi hanya memiliki
        satu label konsensus panel (bukan dua penilaian independen), kesepakatan
        antarpakar tidak dapat dihitung untuk Fase~2.
\end{itemize}

\subsection{Definisi Operasional}\label{subsec:definisi-operasional}

Tabel~\ref{tab:definisi} merangkum definisi operasional istilah-istilah
teknis utama yang digunakan sepanjang penelitian ini.

\begin{table}[htbp]
\centering
\caption{Definisi Operasional Istilah Penelitian}
\label{tab:definisi}
\begin{tabular}{@{}p{0.28\linewidth}p{0.65\linewidth}@{}}
  \toprule
  \textbf{Istilah} & \textbf{Definisi Operasional} \\
  \midrule
  \textit{Knowledge Graph}
    & Struktur data graf dalam Neo4j yang terdiri dari \textit{node} berlabel
      \texttt{Grade/Subject}, \texttt{Chapter}, \texttt{Subtopic},
      \texttt{Concept}, dan \texttt{ConceptTarget}, serta \textit{edge}
      berupa relasi struktural (\texttt{HAS\_CHAPTER},
      \texttt{HAS\_CONCEPT}, dll.) maupun relasi semantik
      dalam-buku dan lintas-buku. \\
  Konsep (\texttt{Concept})
    & Entitas pengetahuan Fisika/Kimia/Biologi yang memiliki atribut
      \texttt{name} dan \texttt{description}. \\
  \texttt{ConceptTarget}
    & Node sementara yang menampung \textit{forward reference}, yaitu konsep
      yang disebut sebagai target relasi namun belum diekstrak secara
      lengkap; dipromosi menjadi \texttt{Concept} pada iterasi
      berikutnya. \\
  Relasi Struktural
    & Hubungan hierarkis dalam graf: \texttt{Grade/Subject} $\to$
      \texttt{Chapter} $\to$ \texttt{Subtopic} $\to$
      \texttt{Concept}. \\
  Relasi Implisit
    & Hubungan antarkonsep yang tidak tersurat dalam teks satu buku,
      diinferensikan melalui kemiripan semantik dan klasifikasi LLM. \\
  KGC (\textit{Knowledge Graph Completion})
    & Proses prediksi \textit{triple} (head, relation, tail) yang
      hilang menggunakan \textit{embedding similarity} dan klasifikasi
      LLM untuk menghasilkan relasi lintas-buku. \\
  \textit{Embedding}
    & Representasi vektor numerik dari deskripsi konsep yang dihasilkan
      model \textit{embedding} (\texttt{gemini-embedding-001}, dimensi:
      3072). \\
  \textit{Cosine Similarity}
    & Ukuran kemiripan antara dua vektor \textit{embedding}, dihitung
      sebagai \textit{dot product} dibagi hasil kali norma kedua vektor.\\
  ANN (\textit{Approximate Nearest Neighbor})
    & Algoritma pencarian tetangga terdekat berbasis HNSW pada Neo4j
      \textit{Vector Index} untuk mempercepat identifikasi kandidat
      relasi lintas-dokumen tanpa perbandingan brute-force
      $O(n^2)$. \\
  \textit{In-Context Learning} (ICL)
    & Strategi pemanfaatan LLM dengan menyertakan demonstrasi struktur
      ontologi di dalam instruksi \textit{prompt} tanpa memerlukan
      \textit{fine-tuning}. \\
  \textit{Confidence} (LLM)
    & Skor dalam rentang $[0,1]$ yang dikeluarkan model klasifikasi
      \texttt{LINTAS\_BUKU} bersama setiap prediksi relasi sebagai ukuran
      tingkat keyakinan model. Skor ini dipakai untuk menyaring relasi:
      relasi diterima bila skornya memenuhi ambang
      (\texttt{BERKAITAN\_DENGAN} $\geq 0{,}85$; tipe lain $\geq 0{,}70$),
      sedangkan kandidat dengan skor $< 0{,}5$ tidak diikutkan dalam
      survei pakar. Karena skor ini belum terkalibrasi, keandalannya
      diperiksa kembali pada Subbab~\ref{subsubsec:false-positive}. \\
\end{tabular}
\end{table}

% ============================================================
\section{Teknik Pengumpulan Data}
% ============================================================

Penelitian ini menggunakan dua jenis data yang telah diuraikan pada
Subbab~\ref{subsec:data-primer} dan~\ref{subsec:data-validasi}: data primer
kurikulum (dimanfaatkan pada Tahap~1--4) dan data validasi pakar (digunakan
pada kedua fase evaluasi). Karena sumber dan cakupan kedua jenis data tersebut
telah dibahas, bagian ini berfokus pada teknik pengumpulan data validasi pakar.

\subsection{Cara Pengumpulan Data}\label{subsec:cara-pengumpulan}

Pengumpulan data validasi dilakukan melalui aplikasi web KG Review App
yang dikembangkan menggunakan Next.js, dipilih karena kemampuannya
menangani navigasi kuesioner berskala besar (lebih dari 100
\textit{triple} per mata pelajaran) yang tidak dapat diakomodasi oleh
Google Forms. Alur pengumpulan data adalah sebagai berikut.

\begin{enumerate}
  \item Peneliti mengunggah hasil \textit{knowledge graph} dalam format
        JSON ke panel admin aplikasi.
  \item Peneliti mempublikasikan mata pelajaran sehingga dapat diakses
        pakar.
  \item Pakar menerima kredensial \textit{login} dari
        peneliti.
  \item Pakar masuk ke aplikasi, memilih mata pelajaran yang
        ditugaskan, dan menilai setiap \textit{triple} menggunakan empat
        kategori penilaian (Subbab~\ref{subsec:validasi-manusia}).
  \item Pakar dapat menambahkan \textit{missing triples} dan
        memberikan komentar per relasi.
  \item Setelah seluruh \textit{triple} dinilai, pakar
        mengonfirmasi penyelesaian \textit{review}.
\end{enumerate}

% ============================================================
\section{Pengolahan Data}\label{sec:pengolahan-data}
% ============================================================

Setelah seluruh pakar menyelesaikan penilaian, data rating
diolah melalui tahapan berikut.

\begin{enumerate}
  \item \textbf{Ekspor dan Agregasi:} Data penilaian diekspor dari Redis
        dalam format JSON per mata pelajaran. Pada Fase~1, setiap
        \textit{triple} dalam-buku memiliki dua rating independen yang
        diagregasi menjadi satu label konsensus berdasarkan kerangka
        evaluasi \textit{hybrid}. Pada Fase~2, setiap \textit{edge}
        \textit{completion} lintas-mapel dinilai melalui diskusi panel
        sehingga langsung menghasilkan satu label konsensus tunggal.

  \item \textbf{Penetapan \textit{Ground Truth}:} Jika dua
        pakar sepakat, label konsensus digunakan langsung.
        Jika tidak sepakat, AI Judge (LLM berbasis konteks buku teks)
        menentukan label akhir, kecuali pada Fase~2 yang
        menggunakan diskusi panel lintas-disiplin tiga pakar (satu wakil
        per mata pelajaran).

  \item \textbf{Perhitungan Metrik Pakar:} Fuzzy Precision, Fuzzy Recall, dan
        F1-Score dihitung per bab per mata pelajaran dengan membandingkan
        label awal pakar dari KG Review App terhadap label final
        pada \textit{gold standard}/konsensus, menggunakan formulasi pada
        Subbab~\ref{subsec:validasi-manusia}.

  \item \textbf{Perhitungan Kesepakatan:} Gwet's AC1 dihitung khusus pada
        Fase~1, mengukur konsistensi \textit{inter-rater} sebelum
        konsolidasi menggunakan label mentah dari masing-masing
        \textit{pakar}. \textit{Edge} \textit{completion} pada Fase~2
        dinilai melalui diskusi panel sehingga hanya menghasilkan satu
        label konsensus dan tidak masuk dalam perhitungan
        \textit{inter-rater agreement}.

  \item \textbf{Kalkulasi Metrik Struktural:} Metrik topologi graf
        (AD, \textit{Density}, \textit{Modularity}, TRR, DC) dihitung
        melalui kueri Cypher deterministik di Neo4j Aura pada kondisi
        pre-KGC dan post-KGC.
\end{enumerate}

% ============================================================
\section{Analisis Data}
% ============================================================

Analisis data dilakukan melalui tiga pendekatan yang saling melengkapi,
dirancang untuk menjawab rumusan masalah penelitian secara komprehensif
dari sisi kuantitatif objektif maupun validitas pedagogis.

\subsection{Analisis Struktural}

Analisis struktural menafsirkan perubahan karakteristik topologi
\textit{Knowledge Graph} untuk mengevaluasi dampak proses KGC terhadap
kualitas interkoneksi antarkonsep. Berdasarkan metrik struktural yang
dihitung pada tahap pengolahan data (Subbab~\ref{sec:pengolahan-data}),
hasil perbandingan kondisi pre-KGC dan post-KGC disajikan dalam bentuk
tabel berpasangan per mata pelajaran. Peningkatan AD dan \textit{Density}
setelah KGC diinterpretasikan sebagai indikator bahwa relasi implisit
yang diprediksi LLM berhasil memperkaya kepadatan jaringan pengetahuan.
Sebaliknya, penurunan \textit{Modularity} yang terlalu drastis dapat
mengindikasikan bahwa batas komunitas antarbab menjadi kabur akibat
relasi lintas-disiplin yang berlebihan, sehingga perlu dikaji lebih
lanjut melalui validasi pakar.

\subsection{Analisis Validasi Pakar}

Analisis ini menafsirkan akurasi dan kelengkapan relasi yang dihasilkan KG
dari metrik validasi pakar yang dihitung pada Subbab~\ref{sec:pengolahan-data}.
Nilai Gwet's AC1 dibaca sebagai keandalan kesepakatan antarpakar pada Fase~1,
sedangkan Fuzzy Precision, Fuzzy Recall, dan F1-Score \parencite{harju_evaluating_2023}
ditafsirkan sebagai kinerja relasi hasil sistem berdasarkan perbandingan antara
label awal pakar dan label final pada \textit{ground truth} konsensus
(Subbab~\ref{subsec:validasi-manusia}).

\subsection{Analisis Komparatif}

Analisis komparatif membandingkan kualitas \textit{Knowledge Graph} di
antara ketiga mata pelajaran untuk mengidentifikasi perbedaan
karakteristik struktur pengetahuan antardisiplin sains dalam konteks
Kurikulum Merdeka.

Perbandingan dilakukan pada dua dimensi. Pertama, dimensi struktural
membandingkan nilai metrik graf (AD, \textit{Density},
\textit{Modularity}, TRR) antarmata pelajaran pada kondisi post-KGC.
Kedua, dimensi validasi membandingkan nilai Fuzzy Precision, Fuzzy Recall, dan F1
antarmata pelajaran untuk mengidentifikasi disiplin mana yang paling
diuntungkan oleh pendekatan LLM-Assisted KGC.

Triangulasi antara metrik struktural dan hasil validasi pakar dilakukan
untuk menguji koherensi keduanya: apabila mata pelajaran dengan AD
tinggi juga memperoleh nilai Fuzzy Precision tinggi dari pakar,
maka ini memperkuat argumen bahwa kepadatan koneksi yang dihasilkan
sistem memang mencerminkan interkoneksi yang valid secara pedagogis,
bukan sekadar artefak komputasional. Hasil analisis disajikan dalam
tabel ringkasan tiga kolom (Fisika, Kimia, Biologi) yang memuat seluruh
metrik secara berdampingan.

\ifSubfilesClassLoaded{\printbibliography[title={Daftar Pustaka}]}{}

\end{document}
